\scanpage{38}
\deleted{sur la courbe ``pincée'' $Y'\deleted{\unsure{=\widehat Y/I|=\widehat Y\amalg_I}}$ (déduite de $\widehat Y$ en contractant $I$ en un pt) en identifiant les fibres \inserted{réduites sur les $y_i\in I$} du faisceau inv. trivial sur $\widehat Y$, \deleted{\unsure{en}} à $k$ par les $\lambda_i$ -- donc entre elles. Alors $u+x$ (pour $x\in Y$) \inserted{\unsure{dans|de} $J_Y^1$} est représenté par le faisceau $\underline O_{\widehat Y}[x]$ (associé au diviseur $x$), muni de ces $\unsure{\overline m}$s identifications.]}

Pour \inserted{que $x+u$} \deleted{qu'il} soit de la forme $y$ $(y\in Y)$ il faut déjà que $u+x$ et $y$ aient $\unsure{\overline m}$ image dans $J_{\widehat Y}^1$, ce qui \deleted{(pour un genre $g\neq1$)} \inserted{\unsure{et|si} comme l'image de $u$ dans $J_{\widehat Y}^0$ est nulle,} implique déjà \inserted{(pour un $g\geq1$)} $x=y$, donc $u=0$. Si $g=0$, donc $\nu=\operatorname{card}I\geq3$ à cause de l'hyp. anabélienne, on \deleted{\unsure{remplace}} prend une partie \inserted{$J$,} de $I$ à deux élts, et on remplace $\widehat Y$ par $\widehat Y\setminus J$ pour conclure encore à $x=y$ donc $u=0$.

\margin{Cependant, on ne sait pas que l'image de $u$ dans $J_{\widehat Y\setminus J}^0$ est nulle!}

[En fait, on a prouvé que si $u\in\operatorname{Ker}(J_Y^0-J_{\widehat Y}^0)$, alors $(Y+u)\cap Y=\varnothing$ -- et par ailleurs, s'il existe un ouvert non vide $U$ de $Y$ tel que $U+u\subset Y$, alors $u$ est dans le Ker précédent (\unsure{du moins si} $g\neq1$). Mais ces arguments ne marchent que pour \uline{$g\neq1$}.]
\margin{$g\geq2$?}

Il vaut mieux, dans le cas général, présenter les choses sous forme homologique. Considérons les deux \deleted{applications} \inserted{morphismes} $U\xrightarrow{i}Y$, \unsure{inclusion}
\margin{37}
